\subsection{Ordinarity already forces prime density zero}

\begin{result}{A.33}{Proved}
{Every ordinary completed pair has support of upper relative prime density zero}
Let \((A,B)\in\mathcal G_{\rm ord}\) be an ordinary completed pair over
\(\mathbb Q\).  No nonabelian section is assumed.  If \(S\) is its
support, then
\[
 \limsup_{X\longrightarrow\infty}
       \frac{\#\{p\in S:p\le X\}}{\pi(X)}=0,
\tag{A.33.1}
\]
where \(\pi(X)\) counts all rational primes up to \(X\).
Thus the density is relative to the primes, not the ordinary integer
density shared by every set of primes.

More precisely, partition the support into the three disjoint sets
\[
 \begin{split}
 S_0&=\{p:v_p(A)>0\},\qquad S_1=\{p:v_p(B)>0\},\\
 S_\infty&=\{p:v_p(A)=v_p(B)<0\},
 \end{split}
\tag{A.33.2}
\]
and put \(C_0=B\), \(C_1=A\), \(C_\infty=-A/B\).
If \(v_{p_i}(C_i)=m_i\ne0\), then for every squarefree positive
integer \(N\) coprime to \(p_i m_i\),
\[
 \limsup_{X\longrightarrow\infty}
       \frac{\#\{p\in S_i:p\le X\}}{\pi(X)}
           \le \frac{\varphi(N)}{N}.
\tag{A.33.3}
\]
If any of these three \(C_i\) has no nonzero valuation, the pair is
one of the rational pairs treated explicitly in the proof.
In particular the density conclusion cannot distinguish an ordinary
ghost from a full-tower survivor: it holds before any survival test.
\end{result}

\begin{proof}
The local unit equation gives exactly the three nonzero valuation
patterns in A.33.2.  Moreover the local representative of \(C_i\)
at a prime in \(S_i\) is a principal unit.  For \(S_0\) and \(S_1\)
these representatives are \(1-t_p\) and \(t_p\), respectively.
For \(S_\infty\) it is
\[
               -\frac{t_p}{1-t_p}
                         =\frac1{1-1/t_p}\in1+p\mathbb Z_p.
\tag{A.33.4}
\]
This uses the actual ordinary local point, not just its valuation.

We prove a finite Kummer assertion for any completed class \(C\)
with one nonzero ordinary integer valuation
\(v_{p_0}(C)=m\).  Fix a squarefree \(N\) with
\(\gcd(N,p_0m)=1\), represent its reduction by
\(c_N\in\mathbb Q^\times\), and choose
\(\alpha_N\in\overline{\mathbb Q}\) with
\(\alpha_N^N=c_N\).  A fixed primitive \(N\)-th root of unity
identifies the actual finite Kummer action with the continuous map
\[
 \begin{split}
 r_{C,N}:G_{\mathbb Q}&\longrightarrow
       (\mathbb Z/N)\rtimes(\mathbb Z/N)^\times,\\
 \sigma&\longmapsto(b_N(\sigma),\chi_N(\sigma)),
 \qquad \sigma(\alpha_N)=\zeta_N^{b_N(\sigma)}\alpha_N .
 \end{split}
\tag{A.33.5}
\]
Its affine action on the roots, labelled by \(z\in\mathbb Z/N\),
is \(z\mapsto\chi_N(\sigma)z+b_N(\sigma)\).
Changing \(\alpha_N\) changes this map by a translation conjugacy;
having a fixed root is independent of that change.

We claim that A.33.5 is onto.  The cyclotomic projection is onto.
At \(p_0\) it is unramified because \(p_0\nmid N\).
On tame inertia the Kummer coordinate is multiplication by
\(v_{p_0}(c_N)\equiv m\pmod N\) on the tame \(N\)-character.
Indeed an \(N\)-th root of the unit part is obtained in an unramified
extension, while the uniformizer part gives that character.
Since \(m\) is a unit modulo \(N\), inertia maps onto the full
translation subgroup \(\mathbb Z/N\).  A subgroup containing all
translations and surjecting to all units is the whole semidirect
product.  This proves the claim without assuming independence of
separately chosen radical fields or passing to an infinite cyclotomic
extension.

Let \(p\nmid N\) be a prime at which \(C\) has a principal-unit
ordinary representative.  Raising to the \(N\)-th power is an
automorphism of \(1+p\mathbb Z_p\), by Hensel's lemma.  Thus the
local Kummer class of \(c_N\) is zero, so one root of
\(X^N-c_N\) belongs to \(\mathbb Q_p\).
All the other roots are obtained by multiplying it by \(\mu_N\),
which is unramified at \(p\).  The finite representation A.33.5 is
therefore unramified at \(p\), and its Frobenius affine permutation
has a fixed point.  These assertions are invariant under the choice
of a prime over \(p\).

For a prime \(\ell\), count the elements of
\(\mathbb F_\ell\rtimes\mathbb F_\ell^\times\) having a fixed
point.  If its slope is different from one, all \(\ell\) translations
have exactly one fixed point.  If its slope is one, only the zero
translation has a fixed point.  The resulting proportion is
\[
 \frac{\ell(\ell-2)+1}{\ell(\ell-1)}
                  =\frac{\ell-1}{\ell}.
\tag{A.33.6}
\]
The Chinese remainder isomorphism identifies the squarefree affine
group with the product of these affine groups, and an element has
a fixed point exactly when each component does.  Its fixed-point
proportion is consequently \(\varphi(N)/N\).
The condition is a union of conjugacy classes.  Chebotarev applied
to the fixed finite splitting field of A.33.5 therefore gives that
relative prime density for the fixed-point set
\citep[Theorem~8.31, printed p.~147]{CohMilneANT}.
Every prime of \(S_i\) outside the finite set dividing \(N\)
belongs to it, by A.33.4.  This proves A.33.3.

Take \(N\) to be the product of the primes up to a bound, omitting
the finitely many divisors of \(p_i m_i\).
The products \(\prod_{\ell\mid N}(1-1/\ell)\) tend to zero.
One elementary proof observes that the reciprocal product over
all primes at most an integer \(M\) includes the harmonic sum
\(\sum_{n\le M}1/n\) in its expansion into positive terms.
Omitting a fixed finite set multiplies the product by a constant.
Thus each \(S_i\) has upper relative prime density zero.
The finite union A.33.2 has the same property.

It remains to justify the possible absence of \(p_i,m_i\).
The power-completion description in E.101 shows that a class with
all prime valuations zero is one of \(1,-1\).
Local Kummer maps \(\mathbb Q_p^\times\to\widehat{\mathbb Q_p^\times}\)
are injective.  This follows directly from the valuation, finite
root-of-unity, and principal-unit factors of \(\mathbb Q_p^\times\):
their intersection of all power subgroups is the identity.
If \(B=1\), ordinarity would give \(t_p=0\), which is excluded;
if \(B=-1\), it gives \(t_p=2\) at every prime.
Similarly \(A=1\) is impossible, and \(A=-1\) gives \(t_p=-1\).
If \(-A/B=1\), the local equation would give
\(-t_p=1-t_p\), which is impossible; if \(-A/B=-1\), it gives
\(t_p=1/2\).  The injectivity of global localization proved in
E.101 then gives, respectively, the rational pairs
\((2,-1)\), \((-1,2)\), and \((1/2,1/2)\).
Their supports are finite.  This covers all exceptional cases.

The finite affine counts are independently enumerated in
\path{code/check_support_density.py}, with saved output and an exact
JSON certificate.  That computation checks the finite combinatorics;
the inertia, Chebotarev, and Euler-product arguments above prove the
density assertion.  No uniform effective support bound follows by
letting \(N\) vary with \(X\): each Chebotarev application here is
for one fixed finite field.  In particular density zero is not
finite support and does not settle A.2 or A.3.
\end{proof}
